\begin{explanationbox}
\textbf{一、先把三个对象分清：切点、垂足、中心垂足}

设切线与椭圆相切于 $P$。切点 $P$ 在椭圆上；$H_1,H_2$ 是从两个焦点向切线作垂线得到的点；中心 $O$ 向切线的垂足另记为 $Q$。\textbf{一般情况下 $P,H_1,H_2,Q$ 不是同一个点。} 本结论说的是 $H_1,H_2$ 在\textbf{主圆}上，不是说它们在椭圆上。

\textbf{二、最容易记住的画面：两焦点在直线同侧}

椭圆及其内部全部位于切线同一侧，因此焦点 $F_1,F_2$ 也位于切线同侧。设切线的向外单位法向量为 $\boldsymbol n=(u,v)$，即 $u^2+v^2=1$，切线写为
\[
ux+vy=h,\qquad h=\sqrt{a^2u^2+b^2v^2}>0.
\]
从焦点往切线作垂线，得到
\[
d_1=h+cu,\qquad d_2=h-cu.
\]
这两段距离并非独立变化：一段变大时另一段必然变小，但增减的数值不必相等；它们相对于同一条切线的中心距对称，其\textbf{平均数}永远等于中心距 $h$，而\textbf{乘积}永远等于 $b^2$：
\[
\underbrace{\frac{d_1+d_2}{2}}_{\text{算术平均}}=h=d_0,
\qquad
\underbrace{\sqrt{d_1d_2}}_{\text{几何平均}}=b.
\]
注意不能仅凭“一个增、一个减”推出积不变；\textbf{真正保证积不变的是切线条件} $h^2=a^2u^2+b^2v^2$ 和 $a^2-c^2=b^2$。

\textbf{三、几何直觉：把一个焦点对切线照镜子}

在切点 $P$ 处利用椭圆的光学反射性质：把 $F_1$ 关于切线 $l$ 对称到点 $M$，则 $F_2,P,M$ 三点共线，且 $P$ 在线段 $F_2M$ 上。因此
\[
F_2M=F_2P+PM=F_2P+PF_1=2a.
\]
由于 $H_1$ 是 $F_1M$ 的中点，$O$ 是 $F_1F_2$ 的中点，\textbf{三角形中位线定理}立即给出
\[
\boxed{OH_1=\tfrac12 F_2M=a}.
\]
所以 $H_1$ 只能在以 $O$ 为圆心、$a$ 为半径的圆上。同理 $H_2$ 也在这个圆上。记忆画面就是：\textbf{“焦点照镜子，连接另一焦点；取中点，得到长轴圆。”}

\textbf{四、为什么距离积是 $b^2$：圆内点幂的解释}

把原切线 $l$ 关于中心 $O$ 对称，得到另一条平行切线 $l'$。从同一焦点 $F_1$ 分别向 $l,l'$ 作垂线，垂足记为 $H,H'$。由于两切线平行，$H,F_1,H'$ 三点在同一直线上；上一性质说明 $H,H'$ 都在主圆上，且因 $F_1$ 在主圆内部，它位于 $H,H'$ 之间。

中心对称把“$F_2$ 到 $l$ 的距离”变成“$F_1$ 到 $l'$ 的距离”，所以
\[
F_1H=d_1,\qquad F_1H'=d_2.
\]
对主圆应用\textbf{相交弦定理（圆内点幂）}：
\[
F_1H\cdot F_1H'=a^2-OF_1^2=a^2-c^2=b^2.
\]
这说明：\textbf{两焦点到同一切线的距离积，本质上是主圆对焦点的点幂绝对值。}

\textbf{五、法向量为什么是最适合本题的语言？}

写成 $y=kx+m$ 会漏掉竖直切线；写成 $Ax+By=D$ 或单位法向量式则没有这个问题。对直线 $ux+vy=h$，任取直线上的两点 $X_1,X_2$，两式相减得
\[
u(x_1-x_2)+v(y_1-y_2)=0.
\]
所以 $\boldsymbol n=(u,v)$ 与直线方向向量 $X_1X_2$ 垂直，即\textbf{法向量就是指向直线正面、用于测量垂直距离的方向}。当 $\boldsymbol n$ 是单位向量时，$h$ 便直接是原点到切线的距离。

\textbf{六、把主圆与原心距联系起来}

利用 $u^2+v^2=1$：
\[
d_0^2=h^2=a^2u^2+b^2v^2=b^2+c^2u^2.
\]
故 $b\le d_0\le a$。切线水平时，$u=0$，$d_0=b$；切线竖直时，$|u|=1$，$d_0=a$。若已知 $d_0$，还可立即求出\textbf{无须确定是哪条切线的}两焦点距离：
\[
\boxed{\{d_1,d_2\}=\{d_0+\sqrt{d_0^2-b^2},\ d_0-\sqrt{d_0^2-b^2}\}}.
\]

\textbf{七、应用识别信号}

题干一出现“\textbf{椭圆切线、两个焦点、两条垂线的长度之积}”，优先用 $d_1d_2=b^2$；出现“\textbf{焦点作垂线的垂足轨迹}”，优先转为主圆 $x^2+y^2=a^2$；出现“\textbf{中心到切线的距离}”，优先用支撑线距离 $d_0$，再与 $d_1+d_2=2d_0$ 连接。\textbf{先判定确实相切，再使用三个结论。}
\end{explanationbox}
